\section{Biểu diễn có thể diễn giải}
\label{sec:ch03-bieu-dien-de-doc}

\index{biểu diễn có thể diễn giải}%
Tính cục bộ chỉ có ích khi các thành phần của lời giải thích có nghĩa với người
đọc. Vì vậy, LIME phân biệt đầu vào gốc $x$ với \tn{biểu diễn có thể diễn
giải}{interpretable representation} $x'$. Bài báo gốc viết $x \in
\mathbb{R}^{d}$ và $x' \in \{0,1\}^{d'}$: $x'$ cho biết thành phần nào đang
có mặt trong biểu diễn có thể diễn giải, còn $x$ vẫn là thứ được đưa vào
$f$~\autocite{p09lime}. Phụ lục~\ref{app:ky-hieu} giữ sự phân biệt này cho các
chương sau.

\index{đặc trưng}%
Với văn bản, một \tn{đặc trưng}{feature} có thể là sự có mặt hay vắng mặt của
một từ. Mô hình thật có thể dùng embedding, vị trí token hoặc những biến đổi
khác mà bảng từ không bộc lộ. Với ảnh, một thành phần của biểu diễn có thể diễn giải
có thể là một vùng điểm ảnh liền kề; mô hình thật lại nhận tensor các kênh màu. Do đó, $x'$
không phải là bản sao đơn giản của $x$. Nó là một quyết định về điều người đọc
được phép gọi tên.

Quyết định này có hậu quả trực tiếp cho \tn{độ trung thực}{faithfulness}. Nếu mô hình phân biệt
thứ tự từ nhưng $x'$ chỉ giữ túi từ, thì mô hình thay thế không thể biểu diễn
tác động của thứ tự ấy. Nếu một super-pixel gộp hai vật thể có vai trò khác
nhau, một trọng số của nó cũng không cho biết vật thể nào đã làm đổi dự đoán.
LIME vẫn có thể trả một lời giải thích ngắn gọn, nhưng độ ngắn gọn đó đến từ
biểu diễn đã bỏ bớt thông tin.

